\section{Quitar Language: por qué Language es veneno}

Esta sección aborda el juicio técnico central del título del episodio. Muchas rutas VLA traducen primero las señales de los sensores a language tokens y después decodifican la trajectory a partir de esos language tokens. Liu Xianming (刘先明) considera que esto genera un cuello de botella, porque la representación lingüística intermedia introduce supervisión humana y compresión semántica humana, lo que perjudica el data scaling. El enfoque de Xiaopeng (小鹏) consiste en quitar esta capa intermedia de language: vision y language entran como entrada y la action se decodifica directamente.

\begin{figure}[H]
\centering
\includegraphics[width=0.94\textwidth]{remove-language-bottleneck.png}
\caption{Quitar el Language Bottleneck: los tokens lingüísticos intermedios limitan el scaling de los datos. Diagrama conceptual propio, elaborado a partir de la conversación de 00:19:00--00:33:53.}
\end{figure}

\begin{knowledgebox}{Lectura de la figura: conservar el lenguaje como entrada y quitarlo como cuello de botella intermedio}
La ruta de la izquierda traduce las señales de los sensores a lenguaje y luego decodifica la trayectoria a partir del lenguaje; la ruta de la derecha hace que Vision+Language entren directamente al modelo y que la Action salga directamente. No se trata de prescindir del lenguaje, sino de no usarlo como señal de supervisión intermedia obligatoria.
\end{knowledgebox}

\subsection{Por qué el lenguaje se convierte en un cuello de botella}

Esta subsección desglosa primero la expresión tajante «veneno». El problema del lenguaje no es que su semántica sea inútil, sino que, en cuanto se vuelve el intermediario obligado de las señales de conducción de bajo nivel, comprime el mundo continuo en etiquetas discretas.

El lenguaje es una abstracción humana de alto nivel, muy adecuada para instrucciones, explicaciones y sentido común; sin embargo, las señales de los sensores en la conducción autónoma contienen una gran cantidad de información continua, geométrica, dinámica y de grano fino. Si primero hay que traducirlas a lenguaje, se pierde mucha información continua y el proceso de entrenamiento se convierte en etiquetado humano o refinement humano. Para la conducción autónoma, que necesita data scaling masivo, esto limita la eficiencia en el uso de los datos.

\begin{importantbox}{El significado preciso de «Language es veneno»}
El lenguaje como entrada es muy valioso; el lenguaje como representación intermedia obligatoria puede ser tóxico. La toxicidad proviene de depender en exceso de la semántica humana, de reducir el aprovechamiento de los datos continuos, de formar un bottleneck y de obstaculizar el entrenamiento autosupervisado o con datos a gran escala.
\end{importantbox}

\subsection{«Quitar la L» y el aprendizaje autosupervisado}

Una de las lecciones importantes del éxito de la IA en el pasado es usar los datos para unsupervised/self-supervised learning. El aprendizaje autosupervisado construye señales de entrenamiento a partir de la propia estructura de los datos y reduce la dependencia del etiquetado humano. Si la conducción autónoma quiere seguir ese camino, debe reducir al mínimo las capas intermedias de supervisión humana y hacer que el modelo aprenda directamente de enormes volúmenes de datos visuales, de sensores y de trayectorias de conducción.

\begin{warningbox}{No hay que entender «quitar el lenguaje» como «prescindir de la capacidad lingüística»}
La conducción autónoma sigue necesitando comprender instrucciones de navegación, la semántica del tránsito y las intenciones humanas. Lo que se quita es el cuello de botella intermedio, no la entrada lingüística en la interacción humano-máquina ni la comprensión semántica dentro del modelo.
\end{warningbox}

\subsection{Resumen de la sección}

«Language es veneno» expresa una cautela frente a las representaciones intermedias. Para que la IA física aproveche plenamente los datos continuos de los sensores, el lenguaje no puede ser el paso estrecho por el que pase toda la información. Un VLA/Action decoding más directo busca mejorar la eficiencia en el uso de los datos y la capacidad de scaling.
